\documentclass[journal]{IEEEtran}
\usepackage{cite}
\usepackage[cmex10]{amsmath}
\usepackage{graphicx}
\usepackage{booktabs}
\usepackage{multirow}
\usepackage{makecell}
\usepackage{hyperref}
\usepackage{orcidlink}
\usepackage{url}
\usepackage{float}
\usepackage{xr}
\externaldocument[SM-]{PS_vs_S2_GRSL_SM}
\input{math_commands.tex}

\title{Beyond Spatial Resolution: Comparing Sentinel-2 and PlanetScope Imagery for Agricultural Plastic Remote Mapping}

\author{Marlon~F.~de Souza\orcidlink{0000-0002-7929-5790},
        Murilo~H.~S.~de~Oliveira\orcidlink{0009-0000-0691-4653},
        Fernando~K.~I.~Fugihara\orcidlink{0009-0009-4749-7361},
        and~Rubens~A.~C.~Lamparelli\orcidlink{0000-0003-4344-1263}% <-this % stops a space


\thanks{Received 25 July 2026; revised dd Month 2026; accepted dd Month
2026. Date of publication 4 November 2024; date of current version
20 Month 2026. This work was supported in part by the Braskem S.A. 
and the S\~ao Paulo Research Foundation - FAPESP (grant no. 2021/05251-8, 2024/06854-6, and 2024/06856-9); 
CNPq (grant no. 305412/2023-0 and 125120/2023-0); and Dow Ilimite-se/Unicamp (grant no. 01-P-416/2024).
(\textit{Corresponding author: Marlon F. de Souza.})}
\thanks{M.F. de Souza is with Luiz de Queiroz College of Agriculture, University of S\~ao Paulo, Piracicaba, SP, 13418 Brazil. e-mail: marlon.souza@usp.br}% <-this % stops a space
\thanks{M.~H.~S.~Oliveira, F.~K.~I.~Fugihara, and~R.~A.~C.~Lamparelli are with Plasticulture Engineering Center, Universidade Estadual de Campinas.}% <-this % stops a space
\thanks{Digital Object Identifier 10.1109/LGRS.2026.xxxxxxx}
}

% The paper headers
\markboth{IEEE Geoscience and Remote Sensing Letters,~Vol.~X, ~2026}{de Souza \MakeLowercase{\textit{et al.}}: Beyond Spatial Resolution: Comparing Sentinel-2 and PlanetScope Imagery for Agricultural Plastic Remote Mapping}

\begin{document}

\maketitle

\begin{abstract}
The choice of remote sensing imagery is critical for balancing accuracy and computational costs in remote sensing applications. This study compares Sentinel-2 (S2) and PlanetScope (PS) for mapping Agricultural Plastic Structures (APS). We hypothesised that PS would outperform S2 and investigated the trade-off between PS's superior spatial and temporal resolution and the increased computational demands associated with its larger data volume, using seven machine learning classifiers. Results showed that S2 imagery with 10 spectral bands consistently outperformed PS, achieving, on average, 2.5\% higher overall accuracy, even though PS image classification took nearly 10x longer. In addition to refuting our hypothesis, these findings reveal a key scalability challenge: the marginal accuracy gains from using more detailed data sources and complex models can become a bottleneck when deploying large-scale mapping. We conclude that S2 imagery is a suitable choice for cost-effective, scalable APS mapping with traditional ML classifiers. Our contribution clarifies certain aspects of APS mapping, including the significance of SWIR bands and the value of runtime as a comparative metric. Future research should investigate deep learning models on this image set comparison.
\end{abstract}

\begin{IEEEkeywords}
agricultural plastic structures, remote sensing, machine learning, classification.
\end{IEEEkeywords}

\IEEEpeerreviewmaketitle

\section{Introduction}
\IEEEPARstart{T}{he} choice of image dataset and machine learning (ML) algorithm directly affects the accuracy and computational cost of land cover (LC) classification \cite{Sheykhmousa2020, Zhang2025, Li2022LCC}. Selecting a suitable and effective imagery source is not straightforward because, in RS image classification, four key resolutions must be considered: spatial, spectral, temporal, and radiometric. Sentinel-2 (S2) and PlanetScope (PS) are among the most widely used collections, each offering distinct advantages and disadvantages. Since 2018, at least 18 studies have compared PS and S2 imagery for LC classification (LCC), so this comparison is neither novel nor groundbreaking in contemporary research \cite{Basheer2024}. However, the results remain inconclusive regarding which collection yields better outcomes.

This study aims to test whether PS's higher spatial and temporal resolution, associated with the recent incorporation of additional spectral bands, may yield improved classification results than S2. We conducted this comparison in a challenging and little-explored case: mapping agricultural plastic structures (APS). We assessed key factors for effective LCC, aiming for a balanced trade-off between accuracy and runtime. Previous studies have discussed little about the scalability required for LCC across extensive geographic areas. This gap underscores the need for an evaluation that accounts for both accuracy and runtime performance in LC mapping workflows. Using large image datasets with many features and high spatial resolution raises concerns about the sustainability of large-scale machine learning applications, which require analysis of processing to enhance performance, reduce computational costs, and minimize environmental impact \cite{deVries2023, Maxwell2018}.

We chose APS as a case study because the increasing presence of plastic in agri-food chains has raised significant environmental concerns due to plastic leakage and microplastics contamination \cite{Ghaffar2022, Veettil2023, Zhang2025}. However, mapping plastic-mulched farmlands (PMF) and plastic greenhouses (PGH) is challenging due to the need for large datasets, clear imagery, and the high spectral similarity between APS and other targets \cite{deSouza2025}. There is demand for global APS maps that include small or sparse PGH and PMF structures \cite{Fabrizi2025, Lu2025, Zhang2025, Tong2024}, yet limited research has examined the suitability of individual satellite imagery sources for APS classification. Although the superior spatial resolution of PS enables the detection of smaller structures, its poorer spectral resolution (due to the absence of certain wavelengths) hampers the detection of plastic materials, creating a dilemma that warrants further examination. Our main contribution is to present results that promote this discussion. We do not include Landsat, the longest-running and most common satellite \cite{Zhang2025, Hansen&Loveland2012}, as a benchmark because its spatial resolution (30 m) is a significant limitation for detecting small or sparse APS.

\section{Material and Methods}
The methodology comprised five main stages: a Systematic Literature Review (SLR), sampling, data preparation, processing, and evaluation. See the Supplementary Material (SM) for further details. We accessed multitemporal S2 imagery via the Google Earth Engine (GEE) API \cite{Google2024} and downloaded PS imagery via the Planet Education and Research Program \cite{PlanetLabs2024}. Then, we generated a 3-time series with a 30-day mean composite for both satellites, following the method presented by \cite{deSouza2025}.

The study area, encompassing two regions of interest (ROIs), is characterized by a high concentration of PMF and PGH in the metropolitan region of Campinas, S\~ao Paulo, Brazil. Each ROI has an area of approximately 330 km\textsuperscript{2} (see SM Figure \ref{SM-fig:ROI}).

The works of \cite{Basheer2024} and \cite{Tayer2023} have suggested testing two different band configurations for each satellite. First, we used all available PS and S2 bands with a spatial resolution of 20 meters or better, yielding eight PS and ten S2 bands. We then compared the satellites using the equivalent four bands: Red (R), Green (G), Blue (B), and Near Infrared (NIR), resulting in visible and NIR (VNIR). This VNIR configuration was intended to provide similar information from both satellite images, leaving the difference solely in spatial resolution and minor sensor variations. We employed pixel-based classification, testing seven ML classifiers selected via the SLR (see section \ref{SM-sec:Model training} in SM) from the Scikit-Learn 1.3.2 Python library \cite{Pedregosa2011}. All classifiers used the default parameters and the same random state (2025).

We selected the most common metrics in the literature (see SM Table \ref{SM-tab:slr_results}), overall accuracy (OA) and Cohen's kappa ($\kappa$), together with training time (TT) and image inference speed (IS). TT and IS denote the time spent fitting the model to the training set and the time required to predict an image using a pre-trained model, respectively. Additionally, we overlaid and visually inspected the most accurate classification results from each satellite (see SM Figure \ref{SM-fig:visual_insp}). In the results, Table \ref{tab_1} presents the impact of replacing the image set from S2 to PS, showing differences in classification metrics (Eq. \ref{eq1}) and the relative changes in TT and IS (Eq. \ref{eq2}). The difference ($\Delta$) for the classification metrics, OA and $\kappa$, is defined in Equation \ref{eq1}.

\begin{equation}
\Delta M = M_{S2} - M_{PS}
\label{eq1}
\end{equation}

Where $M$ represents the metric being compared (either OA or $\kappa$). The relative change in processing time ($t_{\text{rel}}$), applicable to both TT and IS, is calculated as indicated in Equation \ref{eq2}.

\begin{equation}
t_{\text{rel}} = \frac{t_{\text{PS}}-t_{\text{S2}}}{t_{\text{S2}}} = \frac{t_{\text{PS}}}{t_{\text{S2}}}-1
\label{eq2}
\end{equation}

Where $t_{\text{PS}}$ and $t_{\text{S2}}$ represent the processing times for PlanetScope and Sentinel-2, respectively.

\section{Results and Discussion}
Figure \ref{Fig1} provides a summarized analysis of S2 and PS performance in APS mapping across multiple perspectives. By consolidating satellite imagery, algorithms, band configurations, and runtime into a single, exploratory overview, this visualization highlights the overarching trade-off between accuracy and computational cost, especially when selecting the data source. Although interpreting this figure can be complex, it consolidates all the aspects considered in our comparison of PS and S2, enabling a comprehensive analysis. To supplement and individualize the analysis of each attribute, we included in the SM the figures \ref{SM-fig:res_4_all_bands}, which display separate graphs for 4 and all bands; \ref{SM-fig:tt_is_imagery} exhibits IS and TT relative change when switching satellites, and \ref{SM-fig:tt_is_bands} when varying the number of bands; \ref{SM-fig:oa_k} shows OA and $\kappa$ changes based on ML models and imagery sources. The individual SM graphs show mean values with error bars. The trend in metrics for each satellite remained consistent across band configurations, confirming the influence of the source (Fig. \ref{Fig1} and Table \ref{tab_1}).

\begin{figure}[ht]
    \centering
    \includegraphics[width=1\linewidth]{Figures/Figure_1.png}
    \caption{Accuracy and classification time (in logarithmic scale due to the extensive range) of the comparison between satellites, models, and the number of bands.}
    \label{Fig1}
\end{figure}

\begin{table}[ht]
    \centering
    \scriptsize
    \caption{Comparison of PlanetScope and Sentinel-2 performance across seven ML classifiers using four or full bands.}
    \label{tab_1}

    \begin{tabular}{>{\centering\arraybackslash}p{0.8cm} | c | c | c | c | c}
    \toprule
    Imagery & Models & \makecell{$\Delta$OA\\(S2 -- PS)}
& \makecell{$\Delta$$\kappa$\\(S2 -- PS)}
& \makecell{TT relative\\change (\%)}
& \makecell{IS relative\\change (\%)} \\
    \midrule
    \multirow{8}{*}{\makecell[l]{S2 \emph{vs.} PS\\
    {\tiny 4 bands}
    }}
    & CART & 0.0338 & 0.0508 & -44\% & 964\% \\
    & kNN & 0.0169 & 0.0254 & -14\% & 1224\% \\
    & MLP & 0.0139 & 0.0209 & -20\% & 8531\% \\
    & RF & 0.0167 & 0.0250 & -32\% & 930\% \\
    & SGDC & 0.0498 & 0.0747 & 32\% & 938\% \\
    & LinearSVC & 0.0515 & 0.0773 & 2\% & 948\% \\
    & SVC rbf & 0.0111 & 0.0166 & 11\% & 1504\% \\
    \cline{2-6}
    & \makecell{Average \\ 4 bands} & 0.0277 & 0.0415 & -9\% & 2148\% \\
    \midrule
    \multirow{8}{*}{\makecell[l]{S2 \emph{vs.} PS\\
    {\tiny 10 and 8}\\
    {\tiny bands,}\\
    {\tiny respectively}
    }}
    & CART & 0.0293 & 0.0439 & -40\% & 928\% \\
    & kNN & 0.0205 & 0.0307 & -41\% & 1380\% \\
    & MLP & 0.0164 & 0.0246 & 69\% & 989\% \\
    & RF & 0.0200 & 0.0299 & -37\% & 1561\% \\
    & SGDC & 0.0286 & 0.0430 & 39\% & 1199\% \\
    & LinearSVC & 0.0223 & 0.0335 & 20\% & 1056\% \\
    & SVC rbf & 0.0164 & 0.0246 & 61\% & 1759\% \\
    \cline{2-6}
    & \makecell{Average \\ full bands} & 0.0219 & 0.0329 & 10\% & 1267\% \\
    \bottomrule
    \end{tabular}
\end{table}

S2 imagery with 10 spectral bands achieved the highest OA and $\kappa$, and it outperformed PS across all algorithms analyzed (Fig. \ref{Fig1} and Table \ref{tab_1}). On average, S2 had an OA 2.5\% higher and a $\kappa$ 3.7\% higher than PS. Both image datasets showed longer times and improved accuracies for the sets with more bands (Fig. \ref{Fig1}). PS consistently required more time to classify an image, and the runtime also increased with the algorithm's complexity (Table \ref{tab_1}). The size of the image files accounts for the large difference in IS. In the full-band configuration, although PS had fewer bands (8 \emph{vs.} 10 of S2), the total number of pixels (across all bands) was 788\% higher than in S2. The PS image (36.621.865 pixels x 8 bands x 3 times) had a file size four times larger than the S2 image of the same area (3.299.141 pixels x 10 bands x 3 times), and the additional data increased the PS IS by nearly a factor of 10. The satellite did not affect the TT due to the methodological design employed.

We attribute S2's superior performance to contributions from SWIR bands, present in S2 but not in PS \cite{Despini2025}, or to limitations in pixel-based classification with traditional ML classifiers. Non-visible spectral bands, such as SWIR, play a key role in APS mapping \cite{Zhang2025, Veettil2023}, but a definitive conclusion has not yet been reached. It is not entirely clear, but the spectral reflectance presented in SM Figure \ref{SM-fig:reflectance}a shows a significant difference between APS and other common rural LC. The higher spatial resolution of PS increased misclassification, particularly by increasing false positives (FP). In previous LCC studies, \cite{Rahman2020} conducted their work in an urban-dominated area and identified S2 as the most suitable imagery. In contrast, \cite{Acharki2022, Rao2021, Ye2021} focused on ROIs dominated by agricultural and natural fields and concluded that PS was the best choice. However, in a broad literature review, \cite{Sheykhmousa2020} found that higher resolution reduced accuracy in many studies. Our results reinforce that while the worst spatial resolution of S2 may hinder the detection of small structures (false negatives), it could also reduce FP in pixel-based classification. The complex landscape of the study area leads to misclassifications in both satellite imagery, as shown in SM Figure \ref{SM-fig:visual_insp}.

Discussion of computational costs and energy use in ML \cite{GarciaMartin2019} is relevant to the sustainable production of global maps with sparse APS \cite{Lu2025, Zhang2025}. Previous studies have acknowledged SVM as the best algorithm for LC mapping, but none of them have evaluated its classification cost. Figure \ref{Fig1} reaffirm what is already established in the literature: SVC is a computationally intensive algorithm. In our tests, altering the algorithm often affected OA in a way similar to changing the image set.

Over the past 20 years, advances in data availability and sensor integration have improved LCC, including APS remote mapping \cite{Veettil2023}. However, long-term, high-precision mapping remains challenging \cite{Georganos2018, Zhang2025}. Given PS's inferior performance and higher computational demands, S2 is preferable when using traditional ML. The case study was limited to a small area in the Brazilian interior. Scaling to the entire country, with the same computational power, would take 4 days using the fastest classifier (97.5\% OA), and 5850 days (19 years) with the more accurate one (99.0\% OA). The 1.5\% OA gain isn't worth the 1462.5-fold increase in runtime, making it impractical since the APS lifespan is shorter than 19 years.

\subsection{Limitations and future studies}
Although S2 demonstrated superior performance, our methodological design did not incorporate deep learning models or hyperparameter optimization. Deep neural networks (DNNs) have achieved outstanding results in RS problems for some time \cite{Zhang2025}. Given that technologies are continually evolving in parallel, we are currently investigating deep architectures using the most recent satellite imagery as a continuation of this research. Such architectures require tuning to achieve optimal outcomes, which must also be included. We did not test DNNs because they require substantial training data, and we had a limited labeled sample.

Given the high computational costs and energy consumption required to produce accurate global maps with sparse APS, this topic, along with potential alternatives, should be explored further. Therefore, we recommend investigating manifold learning techniques. In many ML applications, attributes do not occupy a large multidimensional space. Instead, the dataset lies on a low-dimensional manifold not apparent in the original space. This is logical for spectral response, where the bands represent small segments of the same underlying variable, and it applies to other aspects of this problem as well. Consequently, studying dimensionality reduction and latent spaces using nonlinear methods appears promising.

\section{Conclusion}
Classifying RS images over large areas is computationally intensive and time-consuming. PS allocated more resources but did not achieve higher accuracy than S2 for APS mapping, refuting our initial hypothesis. We demonstrated the importance of an appropriate combination of datasets, algorithms, and computational resources for performing the target task. In line with the shift toward cost-effective solutions, we incorporated TT and IS as relevant metrics. Among all tested configurations (satellite, bands, and model), S2 imagery with 10 multispectral bands and an MLP proved to be the most balanced in terms of computational cost, accuracy, and consistency.

The importance of SWIR bands for APS mapping should be recognized. Future research could investigate incorporating deep learning into the compared models and applying manifold learning during data preprocessing.

\section*{Code and Data Availability}
The implementation code used in this study is publicly available in the following repository: \url{https://github.com/MuriloHSO/Beyond-Spatial-Resolution-Code}.

\bibliographystyle{IEEEtran}
\bibliography{bibtex/PS_vs_S2_GRSL}

\end{document}
